% =====================================================================
%  Projeto do 2º Bimestre - Sistemas Operacionais (Prof. Jean-Rémi Bourguet)
%  UVV - Engenharia de Software - Turma EW2M - 2026/2
% =====================================================================
\documentclass[12pt, a4paper, oneside, openany, chapter=TITLE, subsection=TITLE, section=TITLE, subsubsection=TITLE, subsubsubsection=TITLE, english, german, brazil]{abntex2}


\usepackage[utf8]{inputenc}
\usepackage[brazilian,hyperpageref]{backref}
\usepackage{pacotes}


\titulo{Conceitos de Sistemas Operacionais e Introdução ao Linux}
\tituloestrangeiro{Operating Systems Concepts and Introduction to Linux}
\autor{Kaike Costa Alves \\ Letícia Storch \\ Mariana Nunes Pignaton \\ Rhayanne Borges \\ Thayrina Cardoso Stange}
\orientador{Prof. Jean-Rémi Bourguet}
\data{2026}
\instituicao{
    Universidade Vila Velha (UVV) \par
    Curso de Engenharia de Software \par
    Sistemas Operacionais -- Turma EW2M
}
\tipotrabalho{Trabalho Acadêmico}
\local{Vila Velha, ES}
\preambulo{Projeto do segundo bimestre apresentado como requisito parcial para aprovação na disciplina de Sistemas Operacionais.}

% informações do PDF
\makeatletter
\hypersetup{
    pdftitle={\@title},
    pdfauthor={Kaike Costa Alves, Letícia Storch, Mariana Nunes Pignaton, Rhayanne Borges, Thayrina Cardoso Stange},
    pdfsubject={\imprimirpreambulo},
    pdfkeywords={sistemas operacionais}{linux}{shell script},
    pdfcreator={LaTeX with abnTeX2},
    colorlinks=true,
    linkcolor=blue,
    citecolor=blue,
    urlcolor=blue
    }
\makeatother
\setlength{\absparsep}{18pt} % ajusta o espaçamento dos parágrafos do resumo

\begin{document}
\pretextual
\imprimircapa
\imprimirfolhaderosto[]*
\makeatother\clearpage

% ---------- Epígrafe ----------
\begin{epigrafe}
    \vspace*{\fill}
    \begin{flushright}
        \textit{``The nice thing about standards is that you have so many to choose from.''}\\
        (Andrew S. Tanenbaum)
    \end{flushright}
\end{epigrafe}
\clearpage

\begin{resumo}[Resumo]
    Este trabalho reúne os sete últimos capítulos da disciplina de Sistemas Operacionais e a
    aplicação prática desses conceitos no sistema GNU/Linux. Na primeira parte são descritos o
    processo de \textit{boot}, o gerenciamento de entrada e saída, a memória secundária e as
    partições, o gerenciamento de arquivos nos sistemas de arquivos, o controle de grupos,
    usuários e permissões, o gerenciamento de processos e o gerenciamento da memória
    primária. Na segunda parte são resolvidos dez exercícios dos capítulos 6 (\textit{Shell}) e 7
    (\textit{Shell Script}) do livro \textit{Introdução ao Linux} \cite{peixoto2013}, com capturas de
    tela do terminal e explicações de cada comando e \textit{script} utilizado. \\
    \vspace{\onelineskip}
    \noindent
    \textbf{Palavras-chave}: sistemas operacionais. Linux. shell script. processos. memória.
\end{resumo}

\begin{resumo}[Abstract]
    \begin{otherlanguage*}{english}
    This work presents the last seven chapters of the Operating Systems course and applies
    these concepts to the GNU/Linux system. The first part describes the boot process,
    input/output management, secondary memory and partitions, file management in file
    systems, groups, users and permissions, process management and primary memory
    management. The second part solves ten exercises from chapters 6 (Shell) and 7 (Shell
    Script) of the book \textit{Introdução ao Linux} \cite{peixoto2013}, with terminal screenshots
    and explanations of each command and script. \\
    \vspace{\onelineskip}
    \noindent
    \textbf{Keywords}: operating systems. Linux. shell script. processes. memory.
    \end{otherlanguage*}
\end{resumo}


\pdfbookmark[0]{\listfigurename}{lof}
\listoffigures*
\clearpage

\pdfbookmark[0]{\listtablename}{lot}
\listoftables*
\clearpage

\pdfbookmark[0]{\lstlistlistingname}{lol}
\renewcommand{\lstlistlistingname}{Lista de códigos}
\lstlistoflistings
\clearpage

\begin{siglas}
    \item[BIOS] \textit{Basic Input/Output System}
    \item[CPU] \textit{Central Processing Unit} (Unidade Central de Processamento)
    \item[DMA] \textit{Direct Memory Access} (Acesso Direto à Memória)
    \item[E/S] Entrada e Saída (\textit{I/O})
    \item[ESP] \textit{EFI System Partition} (Partição de Sistema EFI)
    \item[FAT] \textit{File Allocation Table}
    \item[FIFO] \textit{First In, First Out}
    \item[FS] \textit{File System} (Sistema de Arquivos)
    \item[GID] \textit{Group Identifier} (Identificador de Grupo)
    \item[GPT] \textit{GUID Partition Table}
    \item[GRUB] \textit{GRand Unified Bootloader}
    \item[HD] \textit{Hard Disk} (Disco Rígido)
    \item[LRU] \textit{Least Recently Used}
    \item[MBR] \textit{Master Boot Record}
    \item[MMU] \textit{Memory Management Unit} (Unidade de Gerenciamento de Memória)
    \item[NTFS] \textit{New Technology File System}
    \item[PCB] \textit{Process Control Block} (Bloco de Controle de Processo)
    \item[PID] \textit{Process Identifier} (Identificador de Processo)
    \item[POST] \textit{Power-On Self Test} (Autoteste ao Ligar)
    \item[PPID] \textit{Parent Process Identifier} (Identificador do Processo-Pai)
    \item[RAM] \textit{Random Access Memory}
    \item[RR] \textit{Round-Robin}
    \item[SSD] \textit{Solid State Drive} (Unidade de Estado Sólido)
    \item[TLB] \textit{Translation Lookaside Buffer}
    \item[UEFI] \textit{Unified Extensible Firmware Interface}
    \item[UID] \textit{User Identifier} (Identificador de Usuário)
\end{siglas}

\pdfbookmark[0]{\contentsname}{toc}
\tableofcontents*
\clearpage


\textual % a partir daqui são os elementos textuais.

\chapter*[Introdução]{Introdução}
\addcontentsline{toc}{chapter}{Introdução}

O sistema operacional é a camada de \textit{software} que gerencia o \textit{hardware} e oferece
serviços aos programas \cite{tanenbaum2016}. Este trabalho está organizado em duas partes.
A Parte~\ref{parte:conceitos} descreve os sete capítulos finais da disciplina, do
\autoref{cap:boot} ao \autoref{cap:memoria}. A Parte~\ref{parte:linux} apresenta a resolução
dos exercícios do livro de \citeonline{peixoto2013}, com capturas de tela do terminal.
A divisão das tarefas entre os integrantes está na \autoref{tab:divisao}.

\begin{table}[H]
    \centering
    \caption{Divisão das tarefas entre os integrantes}
    \label{tab:divisao}
    \begin{tabularx}{\linewidth}{lX}
        \toprule
        \textbf{Integrante} & \textbf{Responsabilidade} \\
        \midrule
        Letícia Storch   & Capítulos 1 e 2; Atividades 6.1 e 6.2 \\
        Mariana Nunes Pignaton & Capítulos 3 e 4 \\
        Kaike Costa Alves & Capítulo 5; Atividades 6.3, 6.4, 6.8, 6.9 e 6.10 \\
        Rhayanne Borges  & Capítulos 6 e 7; Atividade 7.6 \\
        Thayrina Cardoso Stange & Atividades 7.9 e 7.11; montagem e revisão do documento \\
        \bottomrule
    \end{tabularx}
    \legend{Fonte: os autores (2026).}
\end{table}

% ---------------------------------------------------------------------
\part{Conceitos de Sistemas Operacionais}\label{parte:conceitos}
% ---------------------------------------------------------------------
% ==== Responsável: Letícia Storch ====
\chapter{Processo de \textit{boot}}\label{cap:boot}

O \textit{boot} (ou inicialização) é o conjunto de etapas que leva o computador do estado
desligado até um sistema operacional pronto para uso. Como a memória principal é volátil, o
\textit{hardware} precisa localizar, em algum dispositivo de armazenamento, o código capaz de
carregar o núcleo (\textit{kernel}) na memória e entregar-lhe o controle \cite{tanenbaum2016}.

\section{Etapas da inicialização}\label{sec:etapas-boot}

Ao ligar a máquina, o processador começa a executar um código gravado em memória não volátil:
o \textbf{\textit{firmware}}, que pode ser a BIOS (modo legado) ou a
UEFI\footnote{Especificações UEFI. Disponível em: \url{https://uefi.org/specifications}. Acesso
em: 19 nov. 2026.}. O \textit{firmware} executa o POST (\textit{Power-On Self Test}), que testa
componentes básicos como memória e teclado, e em seguida procura um dispositivo de
\textit{boot}. A \autoref{fig:boot} resume o fluxo e a \autoref{tab:bios-uefi} compara os dois
tipos de \textit{firmware}.

\begin{figure}[H]
    \centering
    \caption{Fluxo simplificado do processo de \textit{boot} no Linux}
    \label{fig:boot}
    \begin{tikzpicture}[node distance=6mm and 8mm,
        box/.style={rectangle, draw, rounded corners, align=center, fill=gray!10,
                    minimum width=2.6cm, minimum height=1cm, font=\small},
        arr/.style={->, >=stealth, thick}]
        \node[box] (fw) {\textit{Firmware}\\BIOS/UEFI};
        \node[box, right=of fw] (post) {POST};
        \node[box, right=of post] (bl) {\textit{Bootloader}\\(GRUB)};
        \node[box, below=of bl] (k) {\textit{Kernel}\\+ \textit{initramfs}};
        \node[box, left=of k] (init) {\textit{Init}\\(\texttt{systemd})};
        \node[box, left=of init] (user) {Espaço do\\usuário};
        \draw[arr] (fw) -- (post);
        \draw[arr] (post) -- (bl);
        \draw[arr] (bl) -- (k);
        \draw[arr] (k) -- (init);
        \draw[arr] (init) -- (user);
    \end{tikzpicture}
    \legend{Fonte: os autores (2026).}
\end{figure}

\begin{table}[H]
    \centering
    \caption{Comparação entre BIOS e UEFI}
    \label{tab:bios-uefi}
    \begin{tabularx}{\linewidth}{lXX}
        \toprule
        \textbf{Característica} & \textbf{BIOS (legado)} & \textbf{UEFI} \\
        \midrule
        Esquema de partição & MBR & GPT (também aceita MBR) \\
        Tamanho máximo de disco & Cerca de 2\,TiB & Acima de 2\,TiB \\
        Local do \textit{bootloader} & Primeiro setor do disco (MBR, 512 bytes) & Partição de sistema EFI (ESP, em FAT) \\
        Interface de configuração & Texto & Texto ou gráfica \\
        \textit{Secure Boot} & Não & Sim \\
        \bottomrule
    \end{tabularx}
    \legend{Fonte: os autores (2026), com base em \citeonline{silberschatz2015}.}
\end{table}

\section{\textit{Bootloader}}\label{sec:bootloader}

O \textit{bootloader} é o programa responsável por carregar o \textit{kernel}. No Linux, o mais
comum é o GRUB\footnote{Manual do GRUB. Disponível em:
\url{https://www.gnu.org/software/grub/manual/grub/}. Acesso em: 19 nov. 2026.}, que permite
escolher entre os sistemas instalados e passar parâmetros ao \textit{kernel}. Ele carrega também
o \textit{initramfs}, um sistema de arquivos temporário na memória com os \textit{drivers}
necessários para montar o sistema de arquivos raiz.

O menu do GRUB também dá acesso ao \textbf{modo de recuperação} (\textit{recovery mode}) do
Ubuntu, que inicia o sistema com o mínimo de serviços e oferece um terminal de administrador.
Nesse modo, o sistema de arquivos raiz é montado apenas para leitura e precisa ser remontado
com \texttt{mount -o remount,rw /} antes de qualquer alteração, como a troca de uma senha
esquecida com \texttt{passwd}.

\section{\textit{Kernel} e processo \texttt{init}}\label{sec:init}

Com o \textit{kernel} na memória, ele inicializa os dispositivos, monta o sistema de arquivos
raiz e cria o primeiro processo do espaço do usuário (PID~1). Nas distribuições modernas, esse
processo é o \texttt{systemd}, que inicia os serviços em paralelo e define o estado final do
sistema (por exemplo, modo gráfico ou modo texto). Os principais passos, em ordem, são:

\begin{enumerate}
    \item o \textit{firmware} executa o POST e localiza o dispositivo de \textit{boot};
    \item o \textit{bootloader} é carregado e exibe o menu de opções;
    \item o \textit{kernel} e o \textit{initramfs} são carregados na memória;
    \item o \textit{kernel} monta o sistema de arquivos raiz;
    \item o \texttt{init}/\texttt{systemd} inicia os serviços e o ambiente do usuário.
\end{enumerate}

Depois da inicialização, o comando \texttt{systemd-analyze} mostra quanto tempo cada fase do
\textit{boot} levou: o \textit{kernel} e o espaço do usuário e, em máquinas UEFI, também o
\textit{firmware} e o \textit{bootloader}.

% ==== Responsável: Letícia Storch ====
\chapter{Gerenciamento de E/S (I/O)}\label{cap:io}

O gerenciamento de E/S (\textit{I/O}) é a parte do sistema operacional que controla a
comunicação entre o processador e os dispositivos, como discos, teclado, rede e vídeo. Seus
objetivos são oferecer uma interface uniforme para os programas, esconder os detalhes do
\textit{hardware} e coordenar o acesso concorrente \cite{tanenbaum2016}.

\section{Tipos de dispositivos}\label{sec:tipos-dispositivos}

Os dispositivos costumam ser classificados em dois grupos:

\begin{itemize}
    \item \textbf{dispositivos de bloco}: armazenam dados em blocos de tamanho fixo, com
    acesso endereçável (discos e SSDs);
    \item \textbf{dispositivos de caractere}: transferem um fluxo de bytes sem estrutura de
    blocos (teclado, \textit{mouse}, porta serial).
\end{itemize}

No Linux, cada dispositivo é representado por um arquivo especial no diretório
\texttt{/dev}\footnote{Lista de dispositivos do Linux. Disponível em:
\url{https://www.kernel.org/doc/html/latest/admin-guide/devices.html}. Acesso em: 19 nov.
2026.}, o que permite usar as mesmas chamadas de sistema (\texttt{open}, \texttt{read},
\texttt{write}) para arquivos comuns e para \textit{hardware}. No \texttt{ls -l /dev}, a primeira
letra das permissões mostra o tipo: \texttt{b} para dispositivos de bloco (como
\texttt{/dev/sda}) e \texttt{c} para dispositivos de caractere (como \texttt{/dev/tty}).

\section{Camadas do \textit{software} de E/S}\label{sec:camadas-es}

O \textit{software} de E/S é organizado em camadas, de baixo para cima:

\begin{enumerate}
    \item tratadores de interrupção;
    \item \textit{drivers} de dispositivo;
    \item \textit{software} independente de dispositivo (\textit{buffering}, controle de erros,
    alocação);
    \item \textit{software} no espaço do usuário (bibliotecas e programas).
\end{enumerate}

Essa divisão permite trocar o \textit{hardware} sem alterar os programas, porque só o
\textit{driver} conhece os detalhes de cada dispositivo \cite{silberschatz2015}.

\section{Técnicas de transferência de dados}\label{sec:tecnicas-es}

A \autoref{tab:io-tecnicas} apresenta as três técnicas clássicas.

\begin{table}[H]
    \centering
    \caption{Técnicas de transferência de dados em E/S}
    \label{tab:io-tecnicas}
    \begin{tabularx}{\linewidth}{p{3.4cm}XX}
        \toprule
        \textbf{Técnica} & \textbf{Funcionamento} & \textbf{Desvantagem} \\
        \midrule
        E/S programada (\textit{polling}) & A CPU consulta repetidamente o estado do dispositivo. & Desperdiça tempo de CPU. \\
        E/S por interrupção & O dispositivo avisa a CPU ao concluir a operação. & Cada interrupção tem custo de tratamento. \\
        DMA & Um controlador transfere os dados direto entre o dispositivo e a memória. & Exige \textit{hardware} adicional. \\
        \bottomrule
    \end{tabularx}
    \legend{Fonte: os autores (2026), com base em \citeonline{tanenbaum2016}.}
\end{table}

\section{\textit{Buffering}, \textit{cache} e \textit{spooling}}\label{sec:buffering}

O \textit{buffering} usa uma área da memória para compensar a diferença de velocidade entre o
dispositivo e o processo. O \textit{cache} guarda cópias de dados acessados com frequência. O
\textit{spooling} enfileira trabalhos para dispositivos que atendem um processo por vez, como
impressoras.

% ==== Responsável: Mariana Nunes Pignaton ====
\chapter{Memória secundária e partições}\label{cap:particoes}\label{cap:memoria-sec}

A memória secundária é responsável pelo armazenamento permanente de dados e programas.
Diferentemente da memória principal (RAM), ela mantém as informações mesmo com o computador
desligado. Nela ficam o sistema operacional, os aplicativos e os arquivos dos usuários. Em
troca, o acesso é bem mais lento que o da RAM, por isso o sistema operacional precisa
organizar e escalonar o uso dos discos \cite{silberschatz2015}.

\section{Tipos de dispositivos de armazenamento}\label{sec:tipos-armazenamento}

\begin{itemize}
    \item \textbf{Disco rígido (HD):} usa discos magnéticos que giram, com cabeçotes de leitura
    e escrita que se movem até a trilha desejada. É mais lento que o SSD, mas continua sendo
    usado pela grande capacidade e pelo custo menor por \textit{gigabyte}.
    \item \textbf{Unidade de estado sólido (SSD):} usa memória \textit{flash}, sem partes
    mecânicas. Por não ter tempo de busca nem de rotação, é muito mais rápido que o HD, o que
    acelera o \textit{boot} e a abertura de programas.
    \item \textbf{Dispositivos removíveis:} \textit{pendrives}, cartões de memória e HDs
    externos, usados para transportar arquivos e fazer cópias de segurança.
\end{itemize}

\section{Particionamento de discos}\label{sec:particionamento}

O particionamento divide um dispositivo de armazenamento em partes independentes, chamadas
partições. Cada partição pode ter o seu próprio sistema de arquivos e ser usada para guardar
dados ou para instalar um sistema operacional. Isso permite, por exemplo, separar os arquivos
pessoais (\texttt{/home}) dos arquivos do sistema, ou instalar dois sistemas operacionais no
mesmo disco \cite{tanenbaum2016}.

No esquema MBR existem três tipos de partição:

\begin{itemize}
    \item \textbf{primária:} pode guardar um sistema operacional ou arquivos do usuário;
    \item \textbf{estendida:} funciona como um contêiner que permite criar partições lógicas;
    \item \textbf{lógica:} é criada dentro de uma partição estendida.
\end{itemize}

\subsection{Esquemas MBR e GPT}\label{subsec:mbr-gpt}

MBR e GPT são os dois padrões usados para guardar a tabela de partições de um disco. A
\autoref{tab:mbr-gpt} compara os dois.

\begin{table}[H]
    \centering
    \caption{Comparação entre os esquemas de partição MBR e GPT}
    \label{tab:mbr-gpt}
    \begin{tabularx}{\linewidth}{lXX}
        \toprule
        \textbf{Característica} & \textbf{MBR} & \textbf{GPT} \\
        \midrule
        Onde fica a tabela & Primeiro setor do disco (512 bytes) & Início do disco, com cópia de segurança no final \\
        Partições & Até 4 primárias (ou 3 primárias e 1 estendida) & 128 por padrão \\
        Tamanho máximo do disco & Cerca de 2\,TiB & Acima de 2\,TiB \\
        \textit{Firmware} típico & BIOS & UEFI \\
        \bottomrule
    \end{tabularx}
    \legend{Fonte: os autores (2026), com base em \citeonline{tanenbaum2016}.}
\end{table}

O MBR é o padrão mais antigo e tem limites de tamanho de disco e de número de partições. O GPT
é o padrão atual dos computadores com UEFI (ver \autoref{tab:bios-uefi}, no
\autoref{cap:boot}).

\section{Gerenciamento de partições no Linux}\label{sec:particoes-linux}

O Linux tem comandos para ver e administrar os dispositivos de armazenamento e as suas
partições, reunidos no \autoref{lst:particoes}. O \texttt{lsblk}\footnote{Manual do comando
\texttt{lsblk}. Disponível em: \url{https://man7.org/linux/man-pages/man8/lsblk.8.html}.
Acesso em: 20 nov. 2026.} mostra os discos e as partições em forma de árvore (por exemplo,
\texttt{sda}, \texttt{sda1}, \texttt{sda2}). O \texttt{fdisk -l} mostra a tabela de partições
de cada disco. O \texttt{df -h} mostra o espaço usado e livre em cada sistema de arquivos
montado, e o \texttt{du -sh} mostra quanto espaço um diretório ocupa.

\begin{lstlisting}[caption={Comandos para consultar discos e partições no Linux},label={lst:particoes}]
lsblk                  # discos e particoes em arvore
sudo fdisk -l          # tabela de particoes de cada disco
df -h                  # espaco usado e livre em cada sistema de arquivos
du -sh ~/Documentos    # espaco ocupado por um diretorio
\end{lstlisting}

O comando \texttt{df -h} também foi usado na opção 1 do menu da Atividade 7.11
(\autoref{fig:a711a}), em que aparecem a partição \texttt{/dev/sda2}, montada na raiz
(\texttt{/}), e sistemas de arquivos temporários na memória (\texttt{tmpfs}).

% ==== Responsável: Mariana Nunes Pignaton ====
\chapter{Gerenciamento de arquivos em sistemas de arquivos}\label{cap:arquivos}

Um sistema de arquivos (\textit{File System}, FS) organiza, armazena e dá acesso aos dados
gravados em um dispositivo de armazenamento. Ele define como os arquivos são identificados,
organizados em diretórios e recuperados pelo sistema operacional. Sem ele, o disco seria só
uma sequência de blocos sem significado \cite{tanenbaum2016}.

\section{Principais sistemas de arquivos}\label{sec:principais-fs}

Cada sistema de arquivos tem características próprias, resumidas na \autoref{tab:fs}.

\begin{table}[H]
    \centering
    \caption{Comparação entre sistemas de arquivos comuns}
    \label{tab:fs}
    \begin{tabularx}{\linewidth}{lX}
        \toprule
        \textbf{Sistema} & \textbf{Características} \\
        \midrule
        FAT32 & Compatível com quase todos os dispositivos, mas limita cada arquivo a 4\,GB. \\
        exFAT & Feito para \textit{pendrives} e cartões de memória, sem o limite de 4\,GB do FAT32. \\
        NTFS & Padrão do Windows, com permissões, criptografia e suporte a arquivos grandes. \\
        ext4 & Padrão da maioria das distribuições Linux, com \textit{journaling} e bom desempenho. \\
        \bottomrule
    \end{tabularx}
    \legend{Fonte: os autores (2026).}
\end{table}

O \textit{journaling} do ext4\footnote{Documentação do ext4 no \textit{kernel} Linux.
Disponível em: \url{https://docs.kernel.org/filesystems/ext4/}. Acesso em: 20 nov. 2026.}
registra as alterações num diário antes de gravá-las. Se o
computador desligar no meio de uma escrita, o sistema consegue recuperar um estado consistente
sem verificar o disco inteiro \cite{silberschatz2015}.

\section{Organização de diretórios no Linux}\label{sec:diretorios}

No Linux, os arquivos ficam numa árvore de diretórios que começa na raiz, representada por uma
barra (\texttt{/}). Não existem letras de unidade como no Windows: outros discos e partições
aparecem como diretórios dentro dessa mesma árvore. Os principais diretórios são:

\begin{itemize}
    \item \texttt{/home}: arquivos pessoais dos usuários;
    \item \texttt{/etc}: arquivos de configuração do sistema;
    \item \texttt{/bin}: comandos essenciais;
    \item \texttt{/var}: arquivos que mudam com frequência, como os \textit{logs};
    \item \texttt{/tmp}: arquivos temporários;
    \item \texttt{/root}: diretório pessoal do administrador;
    \item \texttt{/dev}: arquivos que representam os dispositivos (ver \autoref{cap:io});
    \item \texttt{/boot}: arquivos usados na inicialização, como o \textit{kernel}
    (ver \autoref{cap:boot}).
\end{itemize}

\section{Operações com arquivos e diretórios}\label{sec:operacoes-arquivos}

O \autoref{lst:arquivos} reúne os comandos básicos para manipular arquivos e diretórios pelo
terminal.

\begin{lstlisting}[caption={Comandos básicos de manipulação de arquivos},label={lst:arquivos}]
ls -la                         # lista tudo, inclusive ocultos, com detalhes
mkdir documentos               # cria um diretorio
touch arquivo.txt              # cria um arquivo vazio
cp arquivo.txt copia.txt       # copia
mv arquivo.txt documentos/     # move (ou renomeia)
rm copia.txt                   # apaga um arquivo
rmdir documentos               # apaga um diretorio vazio
\end{lstlisting}

\section{Caminhos absolutos e relativos}\label{sec:caminhos}

O \textbf{caminho absoluto} indica a localização completa a partir da raiz, como
\texttt{/home/usuario/documentos/arquivo.txt}. O \textbf{caminho relativo} parte do diretório
em que o usuário está: estando em \texttt{/home/usuario}, o mesmo arquivo é
\texttt{documentos/arquivo.txt}. Os atalhos \texttt{.} (diretório atual), \texttt{..}
(diretório pai) e \texttt{\textasciitilde} (diretório pessoal) também são usados nos caminhos.

\section{Montagem de sistemas de arquivos}\label{sec:montagem}

A montagem é o processo de ligar um sistema de arquivos a um diretório da árvore, chamado
ponto de montagem. O comando \texttt{mount} faz a montagem e o \texttt{umount}, a desmontagem:

\begin{lstlisting}[caption={Montagem e desmontagem de uma partição},label={lst:mount}]
sudo mount /dev/sdb1 /mnt      # monta a particao sdb1 em /mnt
sudo umount /mnt               # desmonta
\end{lstlisting}

As partições montadas automaticamente no \textit{boot} ficam listadas no arquivo
\texttt{/etc/fstab}. Além de organizar os dados, o sistema de arquivos também guarda o dono, o
grupo e as permissões de cada arquivo, assunto do \autoref{cap:permissoes}.

% ==== Responsável: Kaike Costa Alves ====
\chapter{Grupos, usuários, poderes e permissões}\label{cap:permissoes}

O Linux é um sistema multiusuário: várias pessoas podem usar o mesmo computador, cada uma com
a sua conta. Cada usuário tem um identificador numérico (UID), pertence a um ou mais grupos
(GID) e recebe permissões específicas para ler, alterar e executar arquivos. Essa separação é
a base da segurança do sistema \cite{silberschatz2015}.

\section{Gerenciamento de usuários}\label{sec:usuarios}

As contas ficam registradas no arquivo \texttt{/etc/passwd}, em que cada linha traz o nome do
usuário, o UID, o GID, o diretório pessoal e o \textit{shell}, separados por dois-pontos. As
senhas criptografadas ficam em \texttt{/etc/shadow}, que só o administrador pode ler. O
\autoref{lst:usuarios} mostra os comandos de criação, alteração e exclusão de contas.

\begin{lstlisting}[caption={Gerenciamento de usuários},label={lst:usuarios}]
sudo useradd -m usuario1           # cria o usuario (com diretorio pessoal)
sudo passwd usuario1               # define a senha
sudo usermod -aG sudo usuario1     # adiciona ao grupo sudo
sudo userdel -r usuario1           # exclui o usuario e o diretorio pessoal
\end{lstlisting}

No exemplo do \texttt{usermod}, a opção \texttt{-aG} acrescenta o usuário a um grupo sem
tirá-lo dos grupos em que já está. Ao entrar no grupo \texttt{sudo}, ele passa a poder usar
privilégios de administrador.

\section{Gerenciamento de grupos}\label{sec:grupos}

Os grupos reúnem usuários para dar permissões a todos de uma vez. Por exemplo, todos os alunos
de uma turma podem ler a pasta de materiais se ela pertencer ao grupo \texttt{estudantes}. Os
grupos ficam em \texttt{/etc/group}.

\begin{lstlisting}[caption={Gerenciamento de grupos},label={lst:grupos}]
sudo groupadd estudantes              # cria o grupo
sudo usermod -aG estudantes usuario1  # coloca o usuario no grupo
groups usuario1                       # mostra os grupos do usuario
sudo groupdel estudantes              # exclui o grupo
\end{lstlisting}

\section{Administrador (\texttt{root}) e o comando \texttt{sudo}}\label{sec:root-sudo}

O administrador do Linux é o usuário \texttt{root} (UID 0), que não tem restrições: pode
instalar programas, mudar configurações, gerenciar usuários e acessar qualquer arquivo. Para
não trabalhar o tempo todo como \texttt{root}, os usuários autorizados usam o
\texttt{sudo}\footnote{Manual do comando \texttt{sudo}. Disponível em:
\url{https://www.sudo.ws/docs/man/sudo.man/}. Acesso em: 20 nov. 2026.}, que executa só aquele
comando com privilégios de administrador e pede a senha do próprio usuário (por exemplo,
\texttt{sudo apt update}). Esse recurso deve ser usado com cuidado, porque um comando
administrativo errado pode comprometer o sistema inteiro.

\section{Permissões de arquivos e diretórios}\label{sec:permissoes-arquivos}

Cada arquivo tem um dono e um grupo, e as permissões são definidas para três categorias:

\begin{itemize}
    \item \textbf{dono} (\texttt{u}): o usuário proprietário do arquivo;
    \item \textbf{grupo} (\texttt{g}): os usuários do grupo do arquivo;
    \item \textbf{outros} (\texttt{o}): todos os demais usuários.
\end{itemize}

Para cada categoria existem três permissões: \textbf{leitura} (\texttt{r}), \textbf{escrita}
(\texttt{w}) e \textbf{execução} (\texttt{x}). Em um diretório, \texttt{r} permite listar o
conteúdo, \texttt{w} permite criar e apagar arquivos dentro dele e \texttt{x} permite entrar
nele com \texttt{cd}.

O \texttt{ls -l} mostra as permissões na primeira coluna, como nesta linha:

\begin{center}
    \texttt{-rw-r-{}-r-{}- 1 usuario1 estudantes 1200 arquivo.txt}
\end{center}

\noindent O primeiro caractere indica o
tipo (\texttt{-} para arquivo comum e \texttt{d} para diretório). Depois vêm três blocos de três
letras: o dono pode ler e escrever (\texttt{rw-}), e o grupo e os outros só podem ler
(\texttt{r-{}-}).

\section{Representação numérica e o comando \texttt{chmod}}\label{sec:chmod}

Cada permissão também tem um valor: leitura vale 4, escrita vale 2 e execução vale 1. A soma dá
um dígito de 0 a 7 para cada categoria, como mostra a \autoref{tab:permissoes}.

\begin{table}[H]
    \centering
    \caption{Exemplos de permissões numéricas no Linux}
    \label{tab:permissoes}
    \begin{tabularx}{\linewidth}{cccX}
        \toprule
        \textbf{Valor} & \textbf{Dono} & \textbf{Grupo e outros} & \textbf{Uso comum} \\
        \midrule
        777 & \texttt{rwx} & \texttt{rwx} / \texttt{rwx} & Acesso total para todos (inseguro). \\
        755 & \texttt{rwx} & \texttt{r-x} / \texttt{r-x} & \textit{Scripts} e diretórios públicos. \\
        644 & \texttt{rw-} & \texttt{r-{}-} / \texttt{r-{}-} & Arquivos de texto comuns. \\
        700 & \texttt{rwx} & \texttt{-{}-{}-} / \texttt{-{}-{}-} & Arquivos privados do dono. \\
        \bottomrule
    \end{tabularx}
    \legend{Fonte: os autores (2026).}
\end{table}

O \texttt{chmod} altera as permissões, tanto em forma numérica quanto com letras
(\autoref{lst:chmod}). Ele foi usado na Atividade 7.9 para tirar e devolver a permissão de
leitura do arquivo \texttt{teste.txt} (\autoref{fig:a79p2}) e em todas as atividades para dar
permissão de execução aos \textit{scripts}.

\begin{lstlisting}[caption={Alteração de permissões, dono e grupo},label={lst:chmod}]
chmod 755 script.sh                     # rwx r-x r-x
chmod 644 documento.txt                 # rw- r-- r--
chmod u+x script.sh                     # da execucao ao dono
chmod g-w documento.txt                 # tira a escrita do grupo
sudo chown usuario1 arquivo.txt         # muda o dono
sudo chown usuario1:estudantes arq.txt  # muda dono e grupo
sudo chgrp estudantes arquivo.txt       # muda so o grupo
\end{lstlisting}

O controle de usuários, grupos e permissões impede que pessoas sem autorização leiam ou alterem
informações. A separação entre usuários comuns e o administrador também reduz o estrago que um
erro ou um programa malicioso pode causar no sistema.

% ==== Responsável: Rhayanne Borges ====
\chapter{Gerenciamento dos processos}\label{cap:processos}

\section{Programa e processo}\label{sec:programa-processo}
O conceito central em qualquer sistema operacional é o de processo. Um programa é um conjunto de instruções guardado em disco, sem consumir recursos. O processo é a instância desse programa em execução, que tem vida mais curta e precisa de memória, dispositivos de E/S e CPU. A \autoref{tab:prog-proc} resume a diferença.

\begin{table}[H]
  \centering
  \caption{Diferença entre programa e processo}\label{tab:prog-proc}
  \begin{tabularx}{\linewidth}{lXX}
    \toprule
    \textbf{Aspecto} & \textbf{Programa} & \textbf{Processo} \\
    \midrule
    Definição & Coleção de instruções que realiza uma tarefa & Instância de um programa em execução \\
    Duração & Mais longa & Mais curta \\
    Local & Armazenado no disco & Carregado na memória principal \\
    Recursos & Não requer recursos & Requer memória, E/S e CPU \\
    \bottomrule
  \end{tabularx}
  \legend{Fonte: adaptado dos slides de \citeonline{bourguet2026processos}.}
\end{table}

\section{Funções de um sistema operacional}

As funções do sistema operacional se dividem em cinco categorias de gerenciamento:

\begin{enumerate}
  \item \textbf{dispositivos de E/S}: aceita a entrada dos dispositivos de entrada e envia o resultado aos de saída, por meio dos \textit{device drivers};
  
  \item \textbf{arquivos}: cria e exclui diretórios e arquivos no sistema de arquivos da memória secundária;
  
  \item \textbf{processos}: gerencia a CPU, criando, agendando e encerrando os processos;
  
  \item \textbf{memória}: controla a alocação e a desalocação dos endereços de memória;
  
  \item \textbf{segurança}: protege o sistema, por exemplo com antivírus e \textit{firewall}.
  
\end{enumerate}

\section{\textit{Time sharing} e vida dos processos}

No \textit{time sharing}, a CPU alterna entre os processos executando cada um por milissegundos. Isso cria o \textbf{pseudoparalelismo}, a ilusão de paralelismo em uma máquina que, na prática, tem uma só CPU. Conceitualmente, cada processo possui sua própria CPU virtual, e o processo passa a maior parte da vida tentando acessar a CPU real.

O processador não sabe qual programa está em execução. Existe apenas um contador de programa físico, e a cada troca acontece o seguinte:
\begin{itemize}
  \item quando um processo é executado, o contador de programa lógico dele é carregado no físico;
  \item quando ele é interrompido, o contador físico é salvo no lógico.
\end{itemize}
Todas as informações do processo interrompido são salvas para que ele continue do ponto onde parou. A troca de um processo por outro é a \textbf{mudança de contexto}, e o conjunto de informações do processo é o seu \textbf{contexto}.

\section{Identificação de processos}

Cada processo é identificado por um número, o PID (\textit{Process Identifier}). Por meio do PID o sistema controla e interfere nos processos. O comando \texttt{pgrep} procura os processos pelo nome e devolve seus PIDs (\autoref{cod:pgrep}).

\begin{lstlisting}[language=bash, caption={Procurando o PID de um processo}, label={cod:pgrep}]
$ pgrep find
11122
\end{lstlisting}

\subsection{O diretório \texttt{/proc}}

O \texttt{/proc} é o sistema de arquivos do \textit{kernel} do GNU/Linux. Ele não contém arquivos reais, apenas informações do sistema, e por isso é um sistema de arquivos virtual. Cada diretório numerado corresponde ao PID de um processo em execução. Dentro de \texttt{/proc/<pid>} os principais arquivos são:

\begin{itemize}
  \item \texttt{cmdline}: argumentos da linha de comando do processo;
  \item \texttt{environ}: variáveis de ambiente visíveis ao processo;
  \item \texttt{fd}: descritores dos arquivos abertos;
  \item \texttt{exe}: link para o executável;
  \item \texttt{status}: estado do processo, uso de memória e outros dados;
  \item \texttt{task/<tid>}: informações das \textit{threads} do processo.
\end{itemize}

\begin{lstlisting}[language=bash, caption={Consultando um processo pelo \texorpdfstring{\texttt{/proc}}{/proc}}, label={cod:proc}]
ls /proc
cd /proc/$(pgrep firefox)
cat cmdline
cat status
cat oom_score
\end{lstlisting}

\section{Bloco de controle e tabela de processos}

O sistema mantém para cada processo um PCB (\textit{Process Control Block}) com todas as suas informações. Os PCBs ficam em uma área exclusiva da memória principal, cujo tamanho é definido pelo número máximo de processos suportados. A \textbf{tabela de processos} guarda um ponteiro para o PCB de cada processo, e o índice da entrada é o próprio PID, o que permite acessar o PCB rapidamente (\autoref{fig:pcb}).

\begin{figure}[H]
  \centering
  \caption{Tabela de processos e blocos de controle (PCB)}\label{fig:pcb}
  \begin{tikzpicture}[font=\small, every node/.style={align=center}]
    \matrix (t) [matrix of nodes, nodes={draw, minimum width=1.2cm, minimum height=0.6cm}, column sep=-\pgflinewidth, row sep=-\pgflinewidth] {
      |[fill=cyan!15]| PID & |[fill=cyan!15]| PCB \\
      1 & $\bullet$ \\
      2 & $\bullet$ \\
      $\vdots$ & $\vdots$ \\
      $n$ & $\bullet$ \\
    };
    \node[draw, rectangle split, rectangle split parts=6, fill=yellow!15, right=3.2cm of t, minimum width=3.4cm] (pcb)
      {Contador de programa \nodepart{two} Registradores \nodepart{three} Estado \nodepart{four} Prioridade \nodepart{five} Espaço de endereçamento \nodepart{six} Arquivos abertos};
    \node[above=0.1cm of pcb] {PCB};
    \draw[-{Stealth}] (t-2-2.center) -- ++(1.2,0) |- (pcb.west);
  \end{tikzpicture}
  \legend{Fonte: adaptado dos slides de \citeonline{bourguet2026processos}.}
\end{figure}

\section{Contexto de um processo}

O PCB guarda o contexto de hardware, o contexto de software e o espaço de endereçamento.

\subsection{Contexto de hardware}

O contexto de hardware armazena o conteúdo dos registradores e contadores da CPU. Esses valores são recarregados antes de o processo retomar a execução. Os principais são mostrados na \autoref{tab:registradores}.

\begin{table}[H]
  \centering
  \caption{Registradores que formam o contexto de hardware}\label{tab:registradores}
  \begin{tabularx}{\linewidth}{llX}
    \toprule
    \textbf{Sigla} & \textbf{Nome} & \textbf{Função} \\
    \midrule
    IR  & \textit{Instruction Register} & Instrução atual que está sendo executada \\
    PC  & \textit{Program Counter} & Endereço da próxima instrução a executar \\
    SP  & \textit{Stack Pointer} & Aponta para o topo da pilha do processo \\
    MAR & \textit{Memory Address Register} & Local dos dados para leitura ou gravação \\
    MBR & \textit{Memory Buffer Register} & Dados transferidos de e para a RAM \\
    AC  & Acumulador & Operações lógicas e aritméticas \\
    PSW & \textit{Program Status Word} & Estado da CPU e do programa em execução \\
    \bottomrule
  \end{tabularx}
  \legend{Fonte: adaptado dos slides de \citeonline{bourguet2026processos}.}
\end{table}

Os dados trafegam entre a CPU e a memória por quatro barramentos: de endereço, de dados, de instruções e de controle.

\subsection{Contexto de software}

O contexto de software reúne as características definidas na criação e na execução do processo:
\begin{itemize}
  \item \textbf{privilégios}: um processo com privilégios elevados acessa recursos críticos, enquanto um com privilégios baixos tem acesso restrito;
  \item \textbf{quotas}: limites de cada recurso que o processo pode alocar, com controle do limite e do uso atual;
  \item \textbf{informações de uso} (\textit{accounting}): tempo de CPU em modo usuário (\textit{user}) e em modo sistema, na execução de funções do \textit{kernel} ligadas ao processo.
\end{itemize}
São exemplos de quotas o número máximo de arquivos abertos, o tamanho máximo de memória principal e secundária alocada, o número máximo de operações de E/S pendentes e o tamanho máximo do \textit{buffer} de E/S\footnote{Disponível em: \url{https://linuxhint.com/set_max_user_processes_linux/}. Acesso em: 19 nov. 2026.}.

\subsection{Espaço de endereçamento}

É a área de memória do processo. Cada processo tem o seu, e o sistema o protege para que os outros processos não o acessem. Esse tema é detalhado na \autoref{cap:memoria}.

\section{Estados de um processo}

Durante a vida do processo, o acesso à CPU e aos recursos se altera, e isso acontece em função de eventos gerados pelo sistema ou pelo próprio processo. Um processo ativo está em um destes estados:

\begin{itemize}
  \item \textbf{execução} (\textit{running}): está usando a CPU;
  \item \textbf{pronto} (\textit{ready}): aguarda sua vez de ser executado;
  \item \textbf{espera} (\textit{waiting}): aguarda um evento ou recurso, como o fim da gravação de um arquivo ou uma determinada hora.
\end{itemize}

O sistema decide o escalonamento, isto é, a ordem e o critério com que os processos em estado pronto acessam a CPU. A \autoref{fig:estados} mostra as transições.

\begin{figure}[H]
  \centering
  \caption{Diagrama de estados de um processo}\label{fig:estados}
  \resizebox{\linewidth}{!}{\begin{tikzpicture}[font=\small, node distance=2.6cm, st/.style={draw, ellipse, minimum width=2cm, minimum height=0.9cm}, >=Stealth]
    \node[st] (cr) {criado};
    \node[st, right=of cr] (pr) {pronto};
    \node[st, right=2.8cm of pr] (ex) {execução};
    \node[st, right=2.2cm of ex] (te) {terminado};
    \node[st, below=2.2cm of $(pr)!0.5!(ex)$] (es) {espera};
    \draw[->] (cr) -- node[above]{admitido} (pr);
    \draw[->] (pr) -- node[above]{despacho} (ex);
    \draw[->] (ex) to[bend right=35] node[above]{interrupção} (pr);
    \draw[->] (ex) -- node[above]{fim} (te);
    \draw[->] (ex) to[bend left=15] node[right, xshift=2mm]{E/S ou evento} (es);
    \draw[->] (es) to[bend left=15] node[left, xshift=-2mm]{fim da E/S} (pr);
  \end{tikzpicture}}
  \legend{Fonte: adaptado dos slides de \citeonline{bourguet2026processos}.}
\end{figure}

\subsection{Criação e eliminação}

Na criação, o sistema adiciona um novo PCB e aloca o espaço de endereçamento na memória. A partir daí ele já reconhece o processo, que fica no estado \textit{criado} até ser admitido na lista de prontos. Na eliminação, os recursos são desalocados e o PCB é removido, mas o sistema ainda mantém na memória algumas informações do processo terminado.

\section{Média de carga}

A média de carga (\textit{load average}) é uma medida do \textit{kernel} do Linux que indica quantas solicitações de recursos do sistema estão pendentes. A cada cinco segundos o sistema conta os processos em espera e calcula uma média móvel exponencial para os últimos 1, 5 e 15 minutos. O comando \texttt{uptime} mostra os valores, e o \texttt{lscpu} informa o número de CPUs, pelo qual a média deve ser dividida.

\begin{lstlisting}[language=bash, caption={Média de carga com \texorpdfstring{\texttt{uptime}}{uptime} e \texorpdfstring{\texttt{lscpu}}{lscpu}}, label={cod:uptime}]
$ uptime
15:29:03 up 14 min, 2 users, load average: 2.92, 4.48, 5.20
$ lscpu | grep "^CPU(s):"
CPU(s):  4
\end{lstlisting}

Com 4 CPUs lógicas, a carga por CPU é $2{,}92/4 = 0{,}73$, $4{,}48/4 = 1{,}12$ e $5{,}20/4 = 1{,}30$ (\autoref{fig:load}). Valores abaixo de 1 indicam uso satisfatório dos recursos, acima de 1 indicam saturação, e uma fila de CPU inativa tem carga 0. Nesse exemplo a carga está diminuindo: no último minuto as CPUs ficaram em uso por cerca de 73\% do tempo, e nos 5 e 15 minutos anteriores o sistema esteve sobrecarregado em torno de 12\% e 30\%.

\begin{figure}[H]
  \centering
  \caption{Média de carga total e por CPU (exemplo do \texorpdfstring{\texttt{uptime}}{uptime})}\label{fig:load}
  \begin{tikzpicture}[x=3.0cm, y=0.65cm, font=\small]
    
    \draw[->] (0.35,0) -- (3.65,0) node[right] {Janela de tempo};
    \draw[->] (0.35,0) -- (0.35,6.3) node[above] {Média de carga};

    \foreach \y in {0,1,2,3,4,5,6} {
      \draw[gray!35] (0.35,\y) -- (3.55,\y);
      \node[left] at (0.35,\y) {\y};
    }

    \fill[blue!60] (1-0.26,0) rectangle (1-0.02,2.92);
    \node[above, font=\scriptsize] at (1-0.14,2.92) {2,92};
    \fill[teal!60] (1+0.02,0) rectangle (1+0.26,0.73);
    \node[above, font=\scriptsize] at (1+0.14,0.73) {0,73};
    \node[below] at (1,0) {1 min};

    \fill[blue!60] (2-0.26,0) rectangle (2-0.02,4.48);
    \node[above, font=\scriptsize] at (2-0.14,4.48) {4,48};
    \fill[teal!60] (2+0.02,0) rectangle (2+0.26,1.12);
    \node[above, font=\scriptsize] at (2+0.14,1.12) {1,12};
    \node[below] at (2,0) {5 min};

    \fill[blue!60] (3-0.26,0) rectangle (3-0.02,5.20);
    \node[above, font=\scriptsize] at (3-0.14,5.20) {5,20};
    \fill[teal!60] (3+0.02,0) rectangle (3+0.26,1.30);
    \node[above, font=\scriptsize] at (3+0.14,1.30) {1,30};
    \node[below] at (3,0) {15 min};

    \fill[blue!60] (0.55,5.65) rectangle (0.75,5.95);
    \node[right, font=\scriptsize] at (0.78,5.8) {Total};

    \fill[teal!60] (1.55,5.65) rectangle (1.75,5.95);
    \node[right, font=\scriptsize] at (1.78,5.8) {Por CPU (total/4)};
  \end{tikzpicture}
  \legend{Fonte: os autores (2026), com os valores dos slides de \citeonline{bourguet2026processos}.}
\end{figure}
% ==== Responsável: Rhayanne Borges ====
\chapter{Gerenciamento da Memória Primária}\label{cap:memoria}
O gerenciamento de memória controla a alocação e a desalocação dos endereços da memória. A memória principal é a que a CPU acessa diretamente para executar um programa, e o sistema precisa dividi-la entre os processos protegendo cada um dos demais \cite{tanenbaum2016}.

\section{Tipos de memória e palavras}
A memória principal é formada por DRAM (\textit{Dynamic RAM}). Nela a carga elétrica vaza lentamente, e por isso um circuito de atualização reescreve os dados periodicamente. Já a SRAM (\textit{Static RAM}), usada na memória \textit{cache} da CPU, não precisa ser atualizada, mas é muito mais cara (\autoref{tab:dram-sram}).

\begin{table}[H]
  \centering
  \caption{Comparação entre DRAM e SRAM}\label{tab:dram-sram}
  \begin{tabularx}{\linewidth}{lXX}
    \toprule
    \textbf{Característica} & \textbf{DRAM} & \textbf{SRAM} \\
    \midrule
    Uso típico & Memória principal (1 a 16 GB) & Memória \textit{cache} (1 a 16 MB) \\
    Atualização & Precisa ser reescrita periodicamente & Não precisa \\
    Custo de produção & Menor & Maior \\
    Transferência & Por blocos & Por palavras \\
    \bottomrule
  \end{tabularx}
  \legend{Fonte: adaptado dos slides de \citeonline{bourguet2026processos}.}
\end{table}

As memórias guardam as informações em sequência, em unidades chamadas palavras. Os registradores guardam uma palavra de dados, de 32 ou 64 bits. A largura da palavra limita a memória endereçável: uma CPU de 32 bits usa no máximo $2^{32}$ endereços, ou 4 GB de RAM, e uma de 64 bits usa até $2^{64}$, cerca de 18,4 EB \cite{machado2013}.

\section{Hierarquia e transferências}
A forma de transferência depende do nível de memória 
(\autoref{fig:hierarquia}):

\begin{itemize}
  \item entre a CPU e a \textit{cache} (SRAM), por \textbf{palavras}, de forma rápida;
  
  \item entre a \textit{cache} e a memória principal (DRAM), por \textbf{blocos}, de forma mais lenta;
  
  \item entre a memória principal e a secundária, por \textbf{páginas} ou segmentos, de forma ainda mais lenta.
  
\end{itemize}

\begin{figure}[H]
  \centering
  \caption{Hierarquia de memória e unidades de transferência}\label{fig:hierarquia}
  \begin{tikzpicture}[font=\small, node distance=1.9cm, bx/.style={draw, minimum width=2.2cm, minimum height=1.2cm, align=center}, >=Stealth]
    \node[bx] (cpu) {CPU};
    \node[bx, right=of cpu] (ca) {\textit{Cache}\\(SRAM)};
    \node[bx, right=of ca] (ram) {Memória\\principal\\(DRAM)};
    \node[bx, right=of ram] (hd) {Memória\\secundária};
    \draw[<->] (cpu) -- node[above]{palavra} node[below]{rápida} (ca);
    \draw[<->] (ca) -- node[above]{bloco} node[below]{lenta} (ram);
    \draw[<->] (ram) -- node[above]{página} node[below]{mais lenta} (hd);
  \end{tikzpicture}
  \legend{Fonte: adaptado dos slides de \citeonline{bourguet2026processos}.}
\end{figure}

\section{\textit{Cache} e \textit{buffer}}

Quando o processador quer ler uma palavra, o sistema segue estes passos:
\begin{enumerate}
  \item testa se a palavra está na memória \textit{cache};
  \item se estiver, ela é entregue direto à CPU;
  \item caso contrário, a palavra é buscada na memória principal.
\end{enumerate}
O \textit{caching} acelera o acesso a dados usados repetidamente, e o \textit{buffering} acelera o fluxo de dados da transmissão. A \textit{cache} tem latência menor que 10 ciclos, contra mais de 100 ciclos da RAM\footnote{Disponível em: \url{https://teachcomputerscience.com/memory/}. Acesso em: 19 nov. 2026.}.

\section{Espaço de endereçamento}

Cada processo tem o seu espaço de endereçamento, protegido pelo sistema contra o acesso dos demais. Ele é organizado, do endereço baixo ao alto, em texto (código binário), dados inicializados, BSS (dados não inicializados), \textit{heap} (memória dinâmica) e pilha (\textit{stack}), com variáveis de ambiente no topo. O \textit{heap} cresce para cima e a pilha cresce para baixo, em direção ao espaço livre entre elas.

O comando \texttt{memstat} mostra o uso de memória pelos processos (\autoref{cod:memstat}), e a \autoref{fig:memstat} mostra os valores do exemplo.

\begin{lstlisting}[language=bash, caption={Uso de memória por processo com \texttt{memstat}}, label={cod:memstat}]
$ sudo apt-get install memstat
$ memstat | sort -nr | head -n 20
1082948k: PID 30047 (/usr/lib/firefox/firefox)
1061960k: PID  2795 (/usr/bin/nautilus)
 991736k: PID 29117 (/usr/bin/evince)
\end{lstlisting}

\begin{figure}[H]
  \centering
  \caption{Uso de memória pelos processos do exemplo do \texttt{memstat}}\label{fig:memstat}
  \begin{tikzpicture}[font=\small]
    % Eixos (0 a 1200 MB -> 0 a 5.0 cm)
    \draw[->] (0,0) -- (8.8,0) node[right] {\small Processo};
    \draw[->] (0,0) -- (0,5.4) node[above] {\small Memória (MB)};

    % Linhas de grade e rótulos do eixo Y
    \foreach \val/\y in {0/0, 300/1.25, 600/2.5, 900/3.75, 1200/5.0} {
      \draw[gray!30] (0,\y) -- (8.4,\y);
      \node[left, font=\scriptsize] at (0,\y) {\val};
    }

    % Barras: {posicao_x/altura_cm/valor_mb/nome}
    \foreach \x/\h/\val/\nome in {
      1/4.4067/1057.6/firefox,
      2/4.3213/1037.1/nautilus,
      3/4.0354/968.5/evince 1,
      4/4.0317/967.6/evince 2,
      5/3.2225/773.4/compiz,
      6/2.5404/609.7/evolution,
      7/2.2263/534.3/texmaker,
      8/2.2058/529.4/zeitgeist
    } {
      \fill[blue!60] (\x-0.26,0) rectangle (\x+0.26,\h);
      \node[above, font=\tiny, rotate=45, anchor=south west] at (\x,\h) {\val};
      \node[below, font=\scriptsize, rotate=40, anchor=north east] at (\x,-0.05) {\nome};
    }
  \end{tikzpicture}
  \legend{Fonte: os autores (2026), com valores do \texttt{memstat} dos slides de \citeonline{bourguet2026processos}, convertidos de kB para MB.}
\end{figure}

\section{\textit{Swapping}}

O número de processos ativos pode superar a capacidade da RAM. Para contornar a falta de recursos usa-se a memória virtual: quando a memória enche, processos em espera são transferidos para o disco, na área de \textit{swap}, e trazidos de volta depois (\textit{swap out} e \textit{swap in}). Qualquer dado a processar precisa voltar para a RAM, pois a área de \textit{swap} é apenas um ponto de espera.

A RAM responde em nanossegundos ($10^{-9}$ s) e o disco em milissegundos ($10^{-3}$ s), uma diferença de cerca de um milhão de vezes. Por isso o uso intenso de \textit{swap} é indesejável e, quando frequente, justifica aumentar a RAM \cite{silberschatz2015}.

\section{Paginação}

Paginação e segmentação são técnicas de alocação de memória não contígua. O endereço lógico é gerado pela CPU durante a execução e é virtual, porque não existe fisicamente: ele serve de referência para acessar a memória física. Como o armazenamento não precisa ser sequencial, a paginação elimina a fragmentação externa.

Na paginação o espaço lógico é dividido em páginas de tamanho fixo, e a memória física em quadros do mesmo tamanho. A tradução usa a \textbf{tabela de páginas}, e cada processo tem um PTBR (\textit{Page-Table Base Register}) que aponta para a sua tabela (\autoref{fig:paginacao}). O processo pode ter as páginas de código carregadas sob demanda: ao iniciar, só a parte inicial é carregada, e as demais vêm quando a execução atinge novas regiões.

Em um endereço de 32 bits dividido em 20 + 12 bits, o número da página usa 20 bits ($2^{20} = 1\,048\,576$ páginas) e o deslocamento (\textit{offset}) usa 12 bits ($2^{12} = 4096$ bytes), de modo que cada página tem 4 KB. No exemplo da \autoref{fig:paginacao}, o endereço virtual com página 2 e \textit{offset} 4095 corresponde ao endereço $2 \times 4096 + 4095 = 12287$. A tabela de páginas leva a página 2 ao quadro 15, e o endereço físico fica $15 \times 4096 + 4095 = 65535$.

\begin{figure}[H]
  \centering
  \caption{Tradução de endereço virtual em físico por paginação}\label{fig:paginacao}
  \begin{tikzpicture}[font=\small, >=Stealth, every node/.style={align=center}]
    \node[draw, minimum width=2.4cm, minimum height=0.8cm, fill=red!10] (vp) {página\\2 (20 bits)};
    \node[draw, minimum width=2cm, minimum height=0.8cm, right=0pt of vp, fill=blue!10] (vo) {\textit{offset}\\4095 (12 bits)};
    \node[above=0.1cm of $(vp.north)!0.5!(vo.north)$] {endereço virtual};
    \node[draw, minimum width=3.2cm, minimum height=1.6cm, below=1.6cm of vp] (pt) {tabela de páginas\\(PTBR)\\página 2 $\to$ quadro 15};
    \node[draw, minimum width=2.4cm, minimum height=0.8cm, right=4.2cm of pt, fill=green!10] (fq) {quadro\\15};
    \node[draw, minimum width=2cm, minimum height=0.8cm, right=0pt of fq, fill=blue!10] (fo) {\textit{offset}\\4095};
    \node[above=0.1cm of $(fq.north)!0.5!(fo.north)$] {endereço físico};
    \draw[->] (vp) -- (pt);
    \draw[->] (pt) -- (fq);
    \draw[->] (vo.east) -- ++(0.5,0) |- (fo.north east) ;
  \end{tikzpicture}
  \legend{Fonte: adaptado dos slides de \citeonline{bourguet2026processos}.}
\end{figure}

\section{Segmentação}

Na segmentação o processo é dividido em partições de tamanho variável, os segmentos. Os detalhes de cada um ficam na tabela de segmentação, com a base (endereço inicial na memória física) e o limite (tamanho). O endereço lógico tem número de segmento e \textit{offset}. Se o \textit{offset} é menor que o limite, o endereço físico é a base mais o \textit{offset}; caso contrário ocorre um erro de endereçamento (\textit{trap}). A \autoref{tab:segmentos} traz a tabela do exemplo.

\begin{table}[H]
  \centering
  \caption{Tabela de segmentos do exemplo}\label{tab:segmentos}
  \begin{tabular}{ccc}
    \toprule
    \textbf{Segmento} & \textbf{Base} & \textbf{Limite} \\
    \midrule
    0 & 500 & 600 \\
    1 & 2500 & 800 \\
    2 & 1500 & 400 \\
    3 & 4600 & 200 \\
    4 & 3800 & 400 \\
    \bottomrule
  \end{tabular}
  \legend{Fonte: adaptado dos slides de \citeonline{bourguet2026processos}.}
\end{table}

Por exemplo, o endereço (segmento 2, \textit{offset} 100) vira $1500 + 100 = 1600$. Já o (segmento 2, \textit{offset} 450) gera erro, pois 450 passa do limite 400.

\begin{table}[H]
  \centering
  \caption{Paginação e segmentação}\label{tab:pag-seg}
  \begin{tabularx}{\linewidth}{lXX}
    \toprule
    \textbf{Aspecto} & \textbf{Paginação} & \textbf{Segmentação} \\
    \midrule
    Tamanho das partições & Fixo (páginas) & Variável (segmentos) \\
    Endereço lógico & Página e \textit{offset} & Segmento e \textit{offset} \\
    Tabela & Tabela de páginas (PTBR) & Tabela de segmentos (base e limite) \\
    Fragmentação externa & Eliminada & Pode ocorrer \\
    \bottomrule
  \end{tabularx}
  \legend{Fonte: os autores (2026), com base em \citeonline{bourguet2026processos} e \citeonline{tanenbaum2016}.}
\end{table}

A relação entre a memória primária e o disco é vista na \autoref{cap:memoria-sec}, e os processos que usam essa memória, na \autoref{cap:processos}.

% ---------------------------------------------------------------------
\part{Introdução ao Linux}\label{parte:linux}
% ---------------------------------------------------------------------
\chapter{\textit{Shell} (Capítulo 6 do livro)}\label{cap:linux6}

Os exercícios a seguir são do capítulo 6 de \citeonline{peixoto2013}. Em cada captura de tela,
o \textit{prompt} mostra o nome do integrante que executou a atividade.

% ==== Atividades 6.1 e 6.2: Letícia Storch ====
\section{Atividade 6.1: exibindo processos em estados específicos}\label{sec:a61}

\textbf{Enunciado:} pesquisar as opções do comando \texttt{ps} e criar um comando que liste
apenas os identificadores e os nomes dos processos que estão no estado R (\textit{running}).
Verificar também o comando \texttt{top}, uma versão interativa do \texttt{ps} que atualiza a
listagem a cada $n$ segundos e ordena os processos por uso de CPU e de memória.

\textbf{Resolução:} o comando \texttt{ps} lista os processos, e a opção \texttt{-e} seleciona
todos eles. Com \texttt{-o}, escolhemos as colunas exibidas: \texttt{pid} (identificador),
\texttt{comm} (nome do comando) e \texttt{stat} (estado). A opção \texttt{-{}-no-headers} tira a
linha de título. O resultado é enviado por um \textit{pipe} (\texttt{|}) ao \texttt{awk}, que
mantém só as linhas cujo terceiro campo começa com a letra \texttt{R} e imprime o PID e o nome
(\autoref{lst:a61}).

\begin{lstlisting}[caption={Atividade 6.1: processos no estado R e \texttt{top}},label={lst:a61}]
ps -eo pid,comm,stat --no-headers | awk '$3 ~ /^R/ {print $1, $2}'
top -d 5
\end{lstlisting}

Na \autoref{fig:a61}, os únicos processos no estado R são o próprio \texttt{ps} e o
\texttt{awk}: eles estão usando a CPU no momento da consulta, enquanto todos os outros estão
dormindo (estado S), esperando algum evento. O \texttt{top -d 5} mostra a mesma informação de
forma contínua, atualizando a tela a cada 5 segundos. No cabeçalho aparecem o total de
processos (26, sendo 1 em execução e 25 dormindo), o uso de CPU e o uso de memória e de
\textit{swap}. Dentro do programa, a tecla \texttt{P} ordena por uso de CPU, \texttt{M} ordena
por uso de memória e \texttt{q} encerra o \texttt{top}.

\printterminal{atividade6_1.jpg}{Atividade 6.1: processos no estado R e uso do \texttt{top}}{fig:a61}

\section{Atividade 6.2: processos em \textit{background}}\label{sec:a62}

\textbf{Enunciado:} a partir do \textit{shell}, executar um programa em \textit{background}.
Observar a árvore hierárquica de processos e as dependências entre eles, observar o PID e o
PPID e comparar com os do \texttt{bash}. Verificar também o comando \texttt{jobs}, que mostra a
situação dos processos em \textit{background} naquele \textit{shell}.

\textbf{Resolução:} colocar o caractere \texttt{\&} no final do comando faz o \textit{shell}
executá-lo em \textit{background} e devolver o \textit{prompt} na hora. Como programa, foi usado
o \texttt{sleep 300}, que só espera 300 segundos, executado duas vezes. Os comandos estão no
\autoref{lst:a62}.

\begin{lstlisting}[caption={Atividade 6.2: processos em \textit{background}},label={lst:a62}]
sleep 300 &            # executa o sleep em background
echo $$                # PID do bash atual
ps -o pid,ppid,comm    # PID e PPID dos processos deste terminal
pstree -p $$           # árvore de processos a partir do bash
jobs -l                # processos em background deste shell
\end{lstlisting}

Ao iniciar o processo, o \textit{shell} mostra o número do \textit{job} entre colchetes e o PID
(\texttt{[2] 1037}). A variável especial \texttt{\$\$} guarda o PID do próprio \textit{shell}:
794. Na saída do \texttt{ps} da \autoref{fig:a62}, os dois \texttt{sleep} (PIDs 979 e 1037) e
o próprio \texttt{ps} têm PPID 794, ou seja, o \texttt{bash} é o processo-pai de todos eles. O
\texttt{pstree -p \$\$} desenha essa mesma dependência como uma árvore, com o \texttt{bash(794)}
na raiz. Por fim, o \texttt{jobs -l} lista os dois \textit{jobs} com número, PID, estado
(\textit{Running}) e comando. O sinal \texttt{+} marca o \textit{job} mais recente, e o
\texttt{-}, o anterior.

\printterminal{atividade6_2.jpg}{Atividade 6.2: PID, PPID, \texttt{pstree} e \texttt{jobs}}{fig:a62}

% ==== Atividades 6.3 e 6.4: Kaike Costa Alves ====
\section{Atividade 6.3: utilizando um \textit{daemon}}\label{sec:a63}

\textbf{Enunciado:} usar o \textit{daemon} \texttt{cron} para fazer, daqui a cinco minutos, um
\textit{backup} dos arquivos \texttt{.log} da pasta \texttt{/var/log}. O arquivo de
\textit{backup} deve ser compactado e colocado dentro da pasta do usuário, usando
\texttt{crontab -e} e \texttt{tar -zcvf destino origem}.

\textbf{Resolução:} um \textit{daemon} é um processo que roda em segundo plano o tempo todo,
sem estar ligado a um terminal. O \texttt{cron} verifica a cada minuto a tabela de tarefas de
cada usuário (a \textit{crontab}) e executa as que estão programadas para aquele horário. Cada
linha da \textit{crontab} tem cinco campos de tempo (minuto, hora, dia do mês, mês e dia da
semana) seguidos do comando.

Para descobrir o horário de daqui a cinco minutos, foi usado o \texttt{date} com a opção
\texttt{-d}, que calcula datas relativas, e o formato \texttt{+\%M \%H}, que imprime só o
minuto e a hora na mesma ordem dos campos da \textit{crontab}. Em seguida, a tarefa foi
adicionada com \texttt{crontab -e} e conferida com \texttt{crontab -l} (\autoref{lst:a63}).

\begin{lstlisting}[caption={Atividade 6.3: tarefa agendada no \texttt{cron}},label={lst:a63}]
date -d "+5 minutes" "+%M %H"      # minuto e hora daqui a 5 minutos
crontab -e                         # edita a tabela de tarefas
crontab -l                         # lista as tarefas agendadas
# linha adicionada na crontab:
06 15 * * * tar -zcvf /home/kaike_alves/backup_logs.tar.gz /var/log/*.log
\end{lstlisting}

A linha \texttt{06 15 * * *} faz o \texttt{cron} executar o \texttt{tar} às 15h06. No
\texttt{tar}, a opção \texttt{-z} compacta com \texttt{gzip}, \texttt{-c} cria o arquivo,
\texttt{-v} mostra os arquivos processados e \texttt{-f} indica o nome do arquivo de destino.
Na última linha da \autoref{fig:a63}, o \texttt{ls -lh} confirma que o arquivo
\texttt{backup\_logs.tar.gz}, com 18\,KB, foi criado às 15:06, exatamente no horário agendado.

\printterminal{atividade6_3.png}{Atividade 6.3: \textit{backup} agendado com o \texttt{cron}}{fig:a63}

\section{Atividade 6.4: usando testes dentro dos \textit{scripts}}\label{sec:a64}

\textbf{Enunciado:} criar um \textit{script} que mostra o usuário atual, o diretório atual e o
\textit{shell} em uso. O livro lembra que o arquivo \texttt{/etc/passwd} lista os usuários e
os seus \textit{shells} e sugere o comando \texttt{cut -d DELIMITADOR -f NUMERO\_DO\_CAMPO}.

\textbf{Resolução:} o usuário atual pode ser obtido com \texttt{whoami} (ou com a variável
\texttt{\$USER}), e o diretório atual com \texttt{pwd}. Para o \textit{shell}, cada linha do
\texttt{/etc/passwd} tem sete campos separados por dois-pontos, e o sétimo é o \textit{shell} de
\textit{login} do usuário. Assim, \texttt{grep \^{}\$USER: /etc/passwd | cut -d: -f7} busca
a linha do usuário e corta só o sétimo campo. Depois de dar permissão de execução com
\texttt{chmod +x}, o \textit{script} foi executado (\autoref{fig:a64}): o usuário é
\texttt{kaike\_alves}, o diretório é \texttt{/home/kaike\_alves} e o \textit{shell} é o
\texttt{/usr/bin/zsh}.

\printterminal{atividade6_4.png}{Atividade 6.4: usuário, diretório e \textit{shell} atuais}{fig:a64}

% ==== Atividades 6.8, 6.9 e 6.10: Kaike Costa Alves ====
\section{Atividade 6.8: utilizando parâmetros e testes}\label{sec:a68}

\textbf{Enunciado:} criar um \textit{script} que verifica se recebeu exatamente um parâmetro na
linha de comando, usando o comando \texttt{test} e saindo com \texttt{exit} se o número de
parâmetros estiver errado. O número de parâmetros fica na variável \texttt{\$\#}.

\textbf{Resolução:} o \textit{script} testa a condição \texttt{test \$\# -ne 1} (ou
\texttt{[ \$\# -ne 1 ]}). Se ela for verdadeira, mostra uma mensagem de uso e termina com
\texttt{exit 1}; se for falsa, mostra o parâmetro recebido, guardado em \texttt{\$1}.

A \autoref{fig:a68} mostra também um erro comum: na primeira tentativa, o \textit{shell}
respondeu \texttt{permission denied}, porque o arquivo criado com o \texttt{nano} ainda não
tinha permissão de execução (\autoref{cap:permissoes}). Depois do \texttt{chmod +x}, o
\textit{script} rodou e mostrou \texttt{Parametro recebido: teste}.

\printterminal{atividade6_8.png}{Atividade 6.8: \textit{script} com teste do número de parâmetros}{fig:a68}

\section{Atividade 6.9: utilizando testes de diretório}\label{sec:a69}

\textbf{Enunciado:} continuar a atividade anterior para verificar se o parâmetro é um
diretório e avisar o usuário se é ou não, usando o teste \texttt{[[ -d \$1 ]]}.

\textbf{Resolução:} o operador \texttt{-d} retorna verdadeiro (código de saída 0, guardado em
\texttt{\$?}) quando o caminho existe e é um diretório. O \textit{script} mantém o teste de
\texttt{\$\#} da Atividade 6.8 e, em seguida, usa um \texttt{if} com \texttt{[[ -d \$1 ]]} para
escolher a mensagem. Na \autoref{fig:a69}, a execução com \texttt{/home} mostra
\texttt{O parametro e um diretorio.}

\printterminal{atividade6_9.png}{Atividade 6.9: teste de diretório com \texttt{-d}}{fig:a69}

\section{Atividade 6.10: listando arquivos passados por parâmetro}\label{sec:a610}

\textbf{Enunciado:} continuar a atividade anterior para receber dois parâmetros, um diretório e
uma letra. Se o primeiro parâmetro for um diretório, o \textit{script} deve listar todos os
arquivos que começam com a letra recebida no segundo parâmetro, usando \texttt{ls \$2*}.

\textbf{Resolução:} agora o teste passa a ser \texttt{[ \$\# -ne 2 ]}. Se \texttt{\$1} for um
diretório, o \textit{script} entra nele e executa \texttt{ls \$2*}. O \textit{shell} expande o
padrão \texttt{\$2*} (por exemplo, \texttt{a*}) para todos os nomes que começam com aquela
letra antes de chamar o \texttt{ls}.

Na \autoref{fig:a610} aparecem dois casos. Com \texttt{/home a}, o \texttt{ls} responde
\texttt{cannot access 'a*'}: não existe nenhum arquivo começando com ``a'' em \texttt{/home}, que
só tem a pasta do usuário. Quando o padrão não encontra nada, o \textit{shell} passa o texto
\texttt{a*} literalmente ao \texttt{ls}, que não acha um arquivo com esse nome. O \texttt{X}
vermelho no \textit{prompt} seguinte indica que o último comando terminou com erro. Com
\texttt{/usr/bin a}, o \textit{script} lista todos os programas desse diretório que começam com
``a'', como \texttt{apt}, \texttt{awk} e \texttt{arch}.

\printterminal{atividade6_10.png}{Atividade 6.10: arquivos de um diretório que começam com uma letra}{fig:a610}

\chapter{\textit{Shell Script} (Capítulo 7 do livro)}\label{cap:linux7}

Um \textit{shell script} é um arquivo de texto com comandos que o interpretador
\textit{Bash} executa em sequência \cite{peixoto2013}. Nesta seção os \textit{scripts} usam
parâmetros (\texttt{\$1}, \texttt{\$2}, \texttt{\$\#}), estruturas condicionais
(\texttt{if}/\texttt{elif}/\texttt{else}), expressões regulares e a estrutura de seleção
\texttt{case}\footnote{Manual oficial do \textit{Bash}. Disponível em:
\url{https://www.gnu.org/software/bash/manual/bash.html}. Acesso em: 13 nov. 2026.}.

Antes de executar um \textit{script} é preciso dar a ele permissão de execução com
\texttt{chmod +x nome.sh}. Depois, ele é chamado com \texttt{./nome.sh}.

% =====================================================================
\section{Atividade 7.6: ainda usando \texttt{grep}}\label{sec:a76}

\textbf{Enunciado:} criar um arquivo de texto chamado \texttt{agenda} com nomes e telefones e desenvolver um \textit{script} que receba um nome como parâmetro e utilize o comando \texttt{grep} para buscar o contato.

O arquivo \texttt{agenda} contém os registros de contatos (um por linha, com nome e telefone), e o \textit{script} \texttt{agenda.sh} (\autoref{lst:agenda}) verifica primeiro se foi passado exatamente um parâmetro (\texttt{\$\# -ne 1}). Em seguida, executa \texttt{grep -i "\$1" agenda} para localizar o nome informado ignorando diferenças entre maiúsculas e minúsculas, exibindo uma mensagem conforme o resultado.

\lstinputlisting[caption={Atividade 7.6 (\texttt{agenda.sh})},label={lst:agenda}]{scripts/agenda.sh}

A execução do \textit{script} no terminal, testando um contato existente e um inexistente, é apresentada na \autoref{fig:a76}.

\printterminal{atividade7_6.png}{Execução da Atividade 7.6}{fig:a76}

% =====================================================================
\section{Atividade 7.9: utilizando expressões regulares}\label{sec:a79}

\subsection{Parte 1: palavra contida em outra}\label{subsec:a79p1}

\textbf{Enunciado:} criar um \textit{script} que recebe duas palavras como parâmetro e verifica
se a primeira está contida na segunda.

\lstinputlisting[caption={Atividade 7.9, parte 1 (\texttt{contida.sh})},label={lst:contida}]{scripts/contida.sh}

O \autoref{lst:contida} primeiro confere, com \texttt{\$\# -ne 2}, se foram passados
exatamente dois parâmetros. Caso contrário, mostra como usar o \textit{script} e encerra com
\texttt{exit 1}. Em seguida, o teste \lstinline[frame=none,numbers=none]{[[ "$2" =~ "$1" ]]} usa o operador
de expressão regular \lstinline[frame=none,numbers=none]{=~} para verificar se o texto de \texttt{\$1} aparece em
qualquer posição de \texttt{\$2}. Como \texttt{\$1} está entre aspas, ele é tratado como texto
literal, e não como padrão. A execução aparece na \autoref{fig:a79p1}.

\printterminal{atividade7_9_parte1.png}{Execução da Atividade 7.9, parte 1}{fig:a79p1}

\subsection{Parte 2: informações de um arquivo texto}\label{subsec:a79p2}

\textbf{Enunciado:} criar um \textit{script} que recebe o nome de um arquivo texto, verifica se o
usuário tem permissão de leitura e, se tiver, mostra o número de caracteres, de palavras e
de linhas usando o comando \texttt{wc}.

\lstinputlisting[caption={Atividade 7.9, parte 2 (\texttt{infoarq.sh})},label={lst:infoarq}]{scripts/infoarq.sh}

No \autoref{lst:infoarq}, o teste \texttt{-e} verifica se o arquivo existe e o teste
\texttt{-r} verifica se o usuário atual pode lê-lo. As opções do \texttt{wc} usadas estão na
\autoref{tab:wc}. A entrada é redirecionada com \texttt{<} para que o \texttt{wc} imprima só o
número, sem repetir o nome do arquivo.

\begin{table}[H]
    \centering
    \caption{Opções do comando \texttt{wc} usadas na Atividade 7.9}
    \label{tab:wc}
    \begin{tabular}{ll}
        \toprule
        \textbf{Opção} & \textbf{O que conta} \\
        \midrule
        \texttt{wc -m} & caracteres \\
        \texttt{wc -w} & palavras \\
        \texttt{wc -l} & linhas \\
        \bottomrule
    \end{tabular}
    \legend{Fonte: adaptado de \citeonline{peixoto2013}.}
\end{table}

Para testar a falta de permissão, a leitura do arquivo foi retirada com
\texttt{chmod -r teste.txt}. Nesse caso o \textit{script} avisa que o usuário não pode ler o
arquivo, como mostra a \autoref{fig:a79p2}.

\printterminal{atividade7_9_parte2.png}{Execução da Atividade 7.9, parte 2}{fig:a79p2}

% =====================================================================
\section{Atividade 7.11: utilizando \texttt{case} com menu de opções}\label{sec:a711}

\textbf{Enunciado:} criar um \textit{script} que mostra um menu com as opções 1)~uso das
partições (\texttt{df -h}); 2)~usuários logados; 3)~data e hora (\texttt{date});
4)~sair. O \textit{script} deve executar o comando escolhido e avisar quando a opção digitada
não existir.

\lstinputlisting[caption={Atividade 7.11 (\texttt{menu.sh})},label={lst:menu}]{scripts/menu.sh}

No \autoref{lst:menu}, o comando \texttt{read} guarda o que o usuário digitou na variável
\texttt{opcao}. O \texttt{case} compara essa variável com cada padrão (\texttt{1)}, \texttt{2)},
\texttt{3)} e \texttt{4)}) e executa apenas o bloco correspondente; cada bloco termina com
\texttt{;;}. O padrão \texttt{*)} funciona como ``qualquer outro valor'' e trata as opções
inválidas. A opção 1 mostra o uso das partições do disco, assunto do
\autoref{cap:particoes}. Na opção 2 foi usado o comando \texttt{w} no lugar do \texttt{who}:
nas versões recentes do Ubuntu, a sessão gráfica (GNOME) não é registrada no arquivo lido pelo
\texttt{who}, que fica vazio. O \texttt{w} mostra o mesmo usuário logado, com o terminal
(\texttt{tty2}), o horário de \textit{login} e o que ele está executando. As
Figuras~\ref{fig:a711a} a~\ref{fig:a711c} mostram cada opção executada.

\printterminal{atividade7_11_opcao1.png}{Atividade 7.11: opção 1 (\texttt{df -h})}{fig:a711a}
\printterminal{atividade7_11_opcao2e3.png}{Atividade 7.11: opções 2 (\texttt{w}) e 3 (\texttt{date})}{fig:a711b}
\printterminal{atividade7_11_invalida_sair.png}{Atividade 7.11: opção inválida e opção 4 (sair)}{fig:a711c}

\chapter*[Conclusão]{Conclusão}
\addcontentsline{toc}{chapter}{Conclusão}

O trabalho mostrou como os conceitos teóricos de sistemas operacionais aparecem no uso
diário do Linux. O processo de \textit{boot}, o gerenciamento de dispositivos, partições,
arquivos, permissões, processos e memória podem ser observados e controlados pelo
terminal. Os \textit{scripts} desenvolvidos na segunda parte mostram que o \textit{shell}
também é uma linguagem de programação, capaz de automatizar tarefas de administração
do sistema com estruturas condicionais, expressões regulares e menus interativos.


\postextual % a partir daqui são os elementos pós-textuais.
\bibliography{referencias}

% ---------- Glossário ----------
\chapter*[Glossário]{Glossário}
\addcontentsline{toc}{chapter}{Glossário}
\begin{description}[style=nextline, leftmargin=1cm]
    \item[\textit{Background}] Execução de um processo em segundo plano, sem bloquear o \textit{shell}; é feita colocando \texttt{\&} no final do comando.
    \item[\textit{Bash}] Interpretador de comandos (\textit{shell}) padrão da maioria das distribuições Linux.
    \item[\textit{Bootloader}] Programa que carrega o \textit{kernel} na memória durante o \textit{boot}, como o GRUB.
    \item[\textit{Case}] Estrutura de seleção do \textit{shell} que compara um valor com vários padrões e executa o bloco correspondente.
    \item[\textit{Cron}] \textit{Daemon} do Linux que executa tarefas agendadas nos horários definidos na \textit{crontab} de cada usuário.
    \item[\textit{Daemon}] Processo que roda em segundo plano o tempo todo, sem estar ligado a um terminal, prestando um serviço ao sistema.
    \item[DMA] Transferência de dados entre o dispositivo e a memória feita por um controlador, sem passar pela CPU.
    \item[\textit{Driver}] \textit{Software} que permite ao sistema operacional controlar um dispositivo específico.
    \item[Escalonador] Parte do sistema operacional que decide qual processo da fila de prontos vai usar a CPU.
    \item[Expressão regular] Padrão de texto formado por caracteres e metacaracteres, usado para buscar ou validar cadeias de caracteres.
    \item[Falta de página] Interrupção gerada quando um processo acessa uma página que não está na memória principal.
    \item[\textit{Journaling}] Técnica em que o sistema de arquivos registra as alterações num diário antes de gravá-las, para se recuperar de desligamentos inesperados.
    \item[Montagem] Ligação de um sistema de arquivos a um diretório da árvore do Linux, chamado ponto de montagem.
    \item[Paginação] Técnica que divide a memória física em quadros e o espaço lógico em páginas do mesmo tamanho, permitindo que um processo ocupe quadros não contíguos.
    \item[Partição] Divisão lógica de um disco, tratada pelo sistema operacional como uma unidade independente.
    \item[Permissão] Direito de leitura (\texttt{r}), escrita (\texttt{w}) ou execução (\texttt{x}) dado ao dono, ao grupo e aos demais usuários de um arquivo.
    \item[POST] Teste automático do \textit{hardware} feito pelo \textit{firmware} logo que o computador é ligado.
    \item[Processo] Programa em execução, com seu próprio contexto (registradores, pilha, memória e estado).
    \item[\textit{Root}] Usuário administrador do Linux (UID 0), sem restrições de acesso.
    \item[\textit{Script}] Arquivo de texto com uma sequência de comandos executados pelo interpretador.
    \item[Sistema de arquivos] Estrutura que organiza os dados de um dispositivo em arquivos e diretórios, como o ext4 e o NTFS.
\end{description}

\end{document}
